%% Đề 18 -- de-kiem-tra-cuoi-hoc-ki-2-toan-9-nam-2023-2024-so-gddt-lang-son.pdf
\deheader{ĐỀ SỐ 18}%
  {Đề kiểm tra Học kỳ — Năm học 2023--2024}%
  {Sở GD\&ĐT Lạng Sơn | Chỉ phần Đại số}
\addcontentsline{toc}{section}{Đề 18 — Đề kiểm tra Học kỳ — Năm học 2023--2024}

\begin{de_cau}{Câu 1}{2,0 điểm}
\begin{enumerate}[label=\alph*)]
\item Giải phương trình: $2x^2 - 7x + 3 = 0$.
\item Rút gọn $Q = \left(\sqrt{x} - \dfrac{x+2}{\sqrt{x}+1}\right) \cdot \dfrac{2}{\sqrt{x}-1}$, $x\geq 0, x\neq 1$.
\end{enumerate}
\end{de_cau}

\begin{loi_giai_cau}
\textbf{a)} $\Delta=49-24=25$. $x=\dfrac{7\pm 5}{4}$: $x=3$ hoặc $x=\dfrac{1}{2}$.

\textbf{b)}
\[\sqrt{x}-\frac{x+2}{\sqrt{x}+1}=\frac{\sqrt{x}(\sqrt{x}+1)-(x+2)}{\sqrt{x}+1}=\frac{x+\sqrt{x}-x-2}{\sqrt{x}+1}=\frac{\sqrt{x}-2}{\sqrt{x}+1}.\]
\[Q=\frac{\sqrt{x}-2}{\sqrt{x}+1}\cdot\frac{2}{\sqrt{x}-1}=\frac{2(\sqrt{x}-2)}{(\sqrt{x}+1)(\sqrt{x}-1)}=\frac{2(\sqrt{x}-2)}{x-1}.\]
\end{loi_giai_cau}

\begin{de_cau}{Câu 2}{2,0 điểm}
\begin{enumerate}[label=\alph*)]
\item Giải hệ phương trình: 
\[ \begin{cases}2x+3y=-1\\x-2y=5\end{cases} \].
\item Tìm tọa độ đỉnh và vẽ parabol $y=-\dfrac{1}{2}x^2$.
\end{enumerate}
\end{de_cau}

\begin{loi_giai_cau}
\textbf{a)} Từ (2): $x=5+2y$. Thế (1): $2(5+2y)+3y=-1\Rightarrow 7y=-11\Rightarrow y=-\frac{11}{7}$.
$x=5+2\cdot(-\frac{11}{7})=\frac{13}{7}$.

\textbf{b)} Đỉnh $O(0;0)$, bề lõm quay xuống. Bảng:
\begin{center}
\begin{tabular}{c|ccccc}$x$&$-2$&$-1$&$0$&$1$&$2$\\\hline$y$&$-2$&$-\frac{1}{2}$&$0$&$-\frac{1}{2}$&$-2$\end{tabular}
\end{center}
\end{loi_giai_cau}

\begin{de_cau}{Câu 3}{2,0 điểm}
Một hồ bơi có thể tích $1800\text{ m}^3$. Có hai máy bơm: nếu máy 1 bơm một mình thì hết 9 giờ, máy 2 bơm
một mình thì hết 12 giờ. Hai máy cùng bơm trong $t$ giờ đầu, sau đó máy 2 hỏng nên chỉ còn máy 1
bơm thêm 2 giờ nữa mới đầy. Tính $t$.
\end{de_cau}

\begin{loi_giai_cau}
Năng suất: máy 1: $\frac{1}{9}$ hồ/giờ, máy 2: $\frac{1}{12}$ hồ/giờ.

$\left(\frac{1}{9}+\frac{1}{12}\right)t + \frac{1}{9}\cdot 2=1$.

$\frac{7}{36}t = 1-\frac{2}{9}=\frac{7}{9}\Rightarrow t=4$ giờ.
\end{loi_giai_cau}

\begin{de_cau}{Câu 4}{1,0 điểm}
Cho phương trình $x^2+2(m-3)x+m^2-3m=0$ (tham số $m$). Tìm $m$ để phương trình có hai nghiệm
dương phân biệt có tổng bằng bình phương hiệu.
\end{de_cau}

\begin{loi_giai_cau}
$\Delta'=(m-3)^2-(m^2-3m)=m^2-6m+9-m^2+3m=-3m+9=3(3-m)$.

Để có 2 nghiệm dương phân biệt:
$\begin{cases}\Delta'>0\Rightarrow m<3\\S=-2(m-3)=2(3-m)>0\Rightarrow m<3\\P=m^2-3m=m(m-3)>0\Rightarrow m<0\text{ hoặc }m>3\end{cases}$

Kết hợp: $m<0$.

Điều kiện $x_1+x_2=(x_1-x_2)^2$: $S=S^2-2P$ (vì $(x_1-x_2)^2=S^2-4P$).

$2(3-m)=[2(3-m)]^2-4m(m-3)=4(3-m)^2+4m(3-m)=4(3-m)[(3-m)+m]=12(3-m)$.

$2=12\Rightarrow$ vô lý. Không có $m$ thỏa mãn với điều kiện đó (kiểm tra lại đề).
\end{loi_giai_cau}
